\section{Dimensión de los patrones y estimación de su suma}
Sobre los sitios de cada línea de las cuatro familias dibujamos un árbol de igualdades,
una arista por igualdad de su coordenada transversal. Una línea con $d$ sitios pierde
su nodo de línea y sus $d$ incidencias, y se reemplaza por $d-1$ aristas. Se conservan
las componentes y los números ciclomáticos del grafo de incidencia. Las restricciones
lineales son las mismas. Trabajamos sólo con patrones realizables con sitios distintos.
En la componente $i$, de $v_i$ sitios y $b_i$ ciclos independientes, escribamos su
rango como $v_i-1+L_i$. Un árbol generador tiene rango $v_i-1$ por eliminación de hojas.

\paragraph{Dos ciclos aún dan dos restricciones independientes.}
Cualquier dependencia entre las ecuaciones asigna coeficientes a sus aristas de modo
que, en cada sitio, la suma con signo de las normales ponderadas es cero.
En el soporte de coeficientes no nulos no hay hojas. En un vértice de grado dos las
normales han de ser paralelas; sus dos aristas están en una misma recta del tablero.
Suprimimos los caminos formados por dichos vértices. Una componente completa de grado
dos sería un ciclo dentro de la misma recta, imposible por el árbol elegido en ella.
Tampoco queda un lazo: el camino recto antes de suprimirlo cerraría el mismo ciclo.
Dos aristas paralelas del grafo reducido darían dos caminos rectos entre los mismos
sitios distintos; éstos determinan una única recta y su unión cerraría un ciclo en ella.
El grafo restante es simple, de grado mínimo tres y tiene al menos cuatro vértices.
Una componente así tiene al menos tres ciclos independientes, pues $2E\ge3V$ da
$E-V+1\ge V/2+1\ge3$. La supresión no altera el número ciclomático.
Un árbol generador con hasta dos cuerdas no puede contener ese soporte. Por tanto
$L_i\ge\min(b_i,2)$, con igualdad si $b_i\le2$. Se permiten dependencias cuando
hay tres ciclos o más; no se añade una dimensión de rango por cada uno de ellos.

\paragraph{Paso a la ley de llegada.}
La dimensión de las colocaciones en las $2v$ coordenadas originales es
\[
 D=2v-\sum_i(v_i-1+L_i)=v+c-\sum_iL_i .
\]
Elegir $D$ coordenadas libres da $O_v(n^D)$ colocaciones. Denotemos por $h$ el número
de índices de fila y columna distintos en el patrón. La prescripción conjunta de llegada
implica una masa retenida $\le C_h(j/n)^h$, sin necesidad de independencia entre ellos.
Con el exponente $p$ de la parte anterior la cota es $C_v n^{D-h}j^{h-p}$.
En los tres casos se tiene $D\le p-1$: la suma de los $L_i$ es cero para bosques,
uno si hay un ciclo, y al menos dos si hay al menos dos ciclos independientes en total.
Si $h\ge p$, usar $j\le n$ concluye. Si $h<p$, una igualdad diagonal entre sitios
distintos es una ecuación no idéntica en las $h$ variables RC: dichos sitios difieren
tanto en fila como en columna. Luego $D\le h-1$ y $j^{h-p}\le1$ da la misma cota.
Así cada patrón aporta en valor absoluto $O_v(1/n)$, y hay sólo un número finito
de patrones (incluidas pendientes y tipos de igualdad) para cada $v$ fijo.

Para $[S]V$ la sustracción de una sola línea se anula salvo si todo el soporte está
en esa línea. En este último caso las demás líneas diagonales son unitarias y el resto
es cero. Los términos de dos conflictos desaparecen para $v\ge5$. Para $v\le4$ se
usan los pesos de dos conflictos y su estimación de media ya revisados en el estudio
interlíneas, parte 1. En particular, dos pares que constituyen todo un soporte
desconectado de cuatro sitios ya fueron sustraídos. Esto da la cota anunciada para $V$.
